%\documentclass[output=paper]{langsci/langscibook}
\ChapterDOI{10.5281/zenodo.20611879}
\documentclass[output=paper]{langscibook} 
\title{Introduction to the book} 
\author{Enrique L. Palancar and Sebastian Fedden}  

\abstract{}

\begin{document}
\maketitle

\noindent If we were to imagine an inflectional morphological system as a galactic universe with its matter and its governing laws, “uninflectedness” would mark the outer boundaries of this system, the “voids", the zones where the gravitational pull of inflection weakens and different structural principles come into play. Viewed in this analogical way, uninflectedness helps trace the limits of inflectional organization and offers valuable insights into how inflection operates and interacts with its periphery. Uninflectedness is thus a complex phenomenon that this book examines in depth through a comprehensive analysis of diverse languages and contexts, aiming to deepen our understanding of inflectional morphology.

In most languages in the world, it is common for the form of words to change to convey grammatical meaning. For example, in English, nouns typically change form to indicate plural (e.g., \emph{men} instead of \emph{man}, \emph{cats} instead of \emph{cat}). However, some nouns, such as SHEEP or DEER, do not change form even in plural contexts (e.g., we cannot say *\emph{these sheeps} or *\emph{these deers}, but \emph{these sheep} or \emph{these deer} are fine, and they mean more than one sheep or deer, respectively). Words like SHEEP and DEER are often referred to as ``invariable''. However, a more technical way to refer to them would be to treat them as ``uninflected'' (i.e., as not inflected) or even as ``uninflecting'', by referring to the fact that they do not inflect. Another way to characterize them is to treat them as ``uninflectable'' a term that highlights their inability to be inflected, that is, manifesting ``uninflectability'' as a grammatical property or characteristic. Authors often use the terms ``uninflectedness'' and ``uninflectability'' interchangeably, because in essence they refer to the same situation. However, one would perhaps like to be more precise and say that ``uninflectability'' refers to the property of being unable to inflect when there is the expectation that they should, while ``uninflectedness'' describes the broader phenomenon by presenting it as lack of inflection, whatever the interpretation one may give as to how or why it happened.

Uninflectedness manifests in various ways across languages and the extent of its presence varies across inflectional systems. In English, uninflected nouns like SHEEP and DEER are somehow anomalies in the English lexicon (i.e., irregular nouns concerning plural formation), but in French, uninflected nouns are the norm, with only a few nouns inflecting for plural (e.g., \emph{oeil/yeux} `eye/eyes', \emph{terminal/terminaux} `terminal/terminals', etc.). Then, just as inflection is sensitive to word class, so is uninflectedness. While some word classes, like adverbs or interjections, commonly fall outside the realm of inflection, the lexicon of a given word class may be divided by uninflectedness. For example, in Basque there are only a handful of verbs which inflect for tense-aspect-mood and the person-number of their arguments, while most inflect periphrastically by auxiliaries (e.g., \emph{da-kar-t} (\textsc{abs3sg}-bring-\textsc{erg1sg}) `I bring it' vs. \emph{erama-ten d-u-t} (take\textsc{-verbal.noun} \textsc{abs3sg}-have-\textsc{erg1sg}) `I take it'). It may also affect entire paradigms or create splits within them, as seen with the noun MUZEUM `museum' in Polish, which is uninflected in the singular, but which inflects normally in the plural.

Uninflectedness can also be observed in agreement. Verbs in Nakh-Dagestanian languages agree with their absolutive argument, but some fail to do so and thus fail to inflect. Systematic uninflectedness also exists. English adjectives whose stem has more than two syllables fail to inflect for comparative or superlative (e.g. \emph{more/most interesting}). Similarly, Spanish adjectives with stems ending in /Xe\#/ do not agree in gender; that is, they are uninflectable for gender (e.g. \emph{la bici grande} `the big bike'/\emph{el coche grande} `the big car' vs. \emph{la bici roja} `the red bike'/\emph{el coche rojo} `the
red car'); raising the question of whether they should be considered as having a different paradigm. In such cases, we can predict the uninflectability of lexemes based on the phonological properties of their stems, suggesting that the phonology limits the lexeme's access to the morphophonological rules of inflection. This would explain why loanwords in many languages pose challenges to speakers inflecting them as native words, and why such words generally fall outside the realm of
morphology and are only inflected once their phonology is adapted to the native system.

All these situations contrast with what Spencer (2020, this volume) calls ``constructional uninflectedness'', where words lose their inflectional properties in specific grammatical constructions. For example, in German, attributive adjectives inflect, but adjectives are uninflectable when used predicatively (e.g. \emph{ein roter Stuhl} `a red.\textsc{nom.sg.m} chair' vs. \emph{dieser Stuhl ist rot} `this chair is red'). Even more strikingly, participial forms of verbs in French do not inflect for the gender of the object argument in perfective periphrastic forms with the auxiliary AVOIR `have', but if reference to
the object is anaphoric, like in headed relative clauses or constructions with fronted NPs, the participle agrees in gender (\emph{e.g. j'ai pris la pomme} /jɛ \textbf{pʁi} la pɔm/ `I've taken the apple', but \emph{la pomme que j'ai prise} /la pɔm kə jɛ \textbf{pʁiz}/ `the apple I've taken' or \emph{la pomme, je l'ai prise} /la pɔm jə lɛ \textbf{pʁiz}/ `the apple, I've taken it').

In this, uninflectedness contributes to understanding partial rule systems \citep{spencer2020} and why languages might allow them at all over general rules. It also opens inquiries into its typological distribution across inflectional systems and its evolutionary dynamics. Given the cognitive load of partial rule systems, one might expect uninflectedness to diminish over time, with items becoming consistently inflected or uninflected, but as we show in the book (see Souag, this volume on Berber) uninflectedness can also be regarded as a stable feature through time in an entire language family. At the opposite end, increasing uninflectedness leads to the death of inflection, as it has happened to case in nouns in some Southern Slavic languages (see Baerman et al., this volume) or in French (see Esher, this volume).

This book explores uninflectedness comprehensively for the first time, assembling a collection of edited articles that examine the phenomenon theoretically and descriptively \emph{with the purpose of building a better typologically-informed theory of inflection}. Detailed analyses of various languages highlight the diversity of uninflectedness, with a focus on Romance and Slavic languages due to the availability of large corpora for investigating the dynamic aspects of the phenomenon.

The volume consists of 13 chapters whose themes we have divided into five main areas:

\medskip

\noindent I. Situating uninflectedness in the theoretical space of inflection:

\medskip

---In his chapter, ``The dog didn't bark, the noun didn't inflect: a typology of significant absences'', Greville G. Corbett argues that uninflectedness is significant because it is both unexpected and noteworthy. This prompts an exploration of why we expect inflection and why its absence matters. Corbett distinguishes uninflectedness from phenomena like syncretism and defectiveness, suggesting it exists on a continuum from fully inflected to uninflected, like the Polish noun \emph{muzeum}, which is uninflected in the \textsc{singular} but inflected in the \textsc{plural} or some Macedonian adjectives which are uninflected for \textsc{gender} but inflected for \textsc{number}. However, in syntax, uninflected items do not always fit seamlessly into expected slots and may be restricted in certain contexts. Uninflectedness varies across criteria, illustrated by the dramatic variations in Slavonic languages and the significant numbers of uninflecting items in Nakh-Dagestanian languages. Corbett calls for refined definitions and mapping of typological possibilities to fully appreciate the complexity and significance of uninflected items.

---Corbett's theoretical stance on uninflectedness is complemented by Andrew Spencer's chapter ``Some concepts and consequences of uninflectedness,'' which raises three questions regarding uninflectability: First, as uninflectable lexemes do not resist derivation, transpositions or evaluative morphology. what counts as `inflection'? Then there is the question about the relationship of uninflectedness to paradigm organization and Spencer asks if uninflectability is not just mass syncretism. The third issue discussed is the relationship to lexical insertion, where Spencer critiques the ``Default Exponence Principle'', which states that, by default, the form of any word undergoing lexical insertion is the lexical root, even in the case of a standardly inflecting lexeme. For uninflectable lexemes he proposes a model of lexical insertion using a coindexing mechanism over pairings of a lexeme's paradigm and syntactic terminals. For uninflecting lexemes, a virtual form/realized paradigm is defined without feature specification. Full constructional uninflectability occurs when syntax lacks expected morphosyntactic features, preventing pairing with the paradigm, thus treating the lexeme as uninflecting. Spencer's chapter highlights the complexity of uninflectedness, suggesting that it involves more than a simple absence of inflection,
affecting derivation, transpositions, and evaluative morphology, and thus challenges existing morphological theories.

\medskip

\noindent II. Conditions for uninflectedness:

\medskip

---Ursula Doleschal identifies the lexeme classes most frequently subject to uninflectedness across several Slavic languages, including Czech, Polish, Russian, Serbian, Slovene, and Slovak. These classes predominantly comprise proper names, loanwords, and abbreviations, and most of her examples are nouns. The goal of Doleschal's study is to provide a comprehensive understanding of the factors influencing uninflectedness in Slavic noun classes by highlighting the varying degrees of integration of nouns within the inflectional systems of the different languages under study. Illustrating that each word class exhibits a continuum from a high degree of uninflectedness to nearly complete inflection, Doleschal argues that the absence of inflectional integration in nouns relates to essential features of declension classes, such as stem-ending, base form, and grammatical \textsc{gender}. These features can vary in restrictiveness, leading to differences among languages. For example, Russian displays a higher degree of uninflectedness in nouns compared to Slovene, but classes of \textsc{feminine} nouns tend to be more restrictive regarding stem structure and base form than those of \textsc{masculine} nouns.

---Then Michele Loporcaro, in the chapter titled ``Uninflectedness as a factor in agreement loss'' argues that some of the Romance languages have gone a long way towards generalizing an isolating word structure, making uninflectedness the rule rather than the exception, at least in some areas of their grammars. The most widely known case is French nominal morphology, but a large subset of the Italo-Romance varieties -- both north and south of Tuscany, whose dialects provided the basis for the standard language -- parallel French in having reduced distinctions in inflection in general, in such a way that agreement has come to be signalled much less systematically than it used to be in Latin and still is in standard Italian. Elaborating on previous work, Loporcaro shows in which ways uninflectedness makes its way into these Romance varieties and in which ways it comes to interact with agreement. The main focus of Loporcaro's study is on the agreement by participles with the direct object, for which he shows that most Romance languages display four-cell paradigms which may show syncretisms or uninflectability, usually as a product of regular sound change.

\medskip

\noindent III. Emergence of uninflectedness:

\medskip

---In their co-authored Chapter, Matthew Baerman, Greville G. Corbett, Alexander Krasovitsky and Maria Kyuseva argue that ongoing loss of case in Serbian and Bulgarian dialects has caused the system to split into competing co-grammars, one where nouns inflect and one where they do not. This split between results in partial rules of different nature. The authors argue that three kinds of split are found, depending on the dialect (as illustrated from Bulgarian). Firstly, there is a split between morphological classes: while some retain case distinctions, others are uninflected. Secondly, classes that retain inflection are nevertheless affected by constructional uninflectedness. Thirdly, an otherwise inflecting morphological class may split on a semantic basis such as human vs. non-human. Asymmetry in morphological change leads to the rise of partial synchronic rules: morphological, syntactic or semantic. These rules persist over time, and change when a historical process moves to a new phase (for example, if uninflectedness spreads to a group within an inflected morphological class, e.g. to inanimates, splitting this class on a semantic basis). By examining these rules in contemporary dialects, and by arranging them chronologically, Baerman et al. show that one can uncover the fine-grained details of this historical process.

Louise Esher's chapter, ``Loss of inflection in the diachrony of French nouns", is a case study that contributes a diachronic perspective on the development of uninflectedness and its interaction with other non-canonical phenomena. Esher argues that Modern French nouns have a ``content paradigm''---the inventory of morphosyntactic feature sets required by syntax---with two cells, corresponding to the values \textsc{singular} and \textsc{plural} of the feature \textsc{number}; this contrast is discernable through agreement patterns. Some nouns such as JOURNAL `newspaper', have two distinct forms (\emph{journal} vs. \emph{journaux}) in the ``realized paradigm''---the array of inflectional wordforms. However, nouns of the majority inflectional class, such as LIVRE `book', display syncretism for \textsc{number}; as the realised paradigm has only a single form (\emph{livre}(\emph{s})), such nouns may also be considered uninflectable. Esher shows that these patterns result from progressive ``loss of inflection'': Medieval French nouns had a content paradigm of four cells, with two values each for the features \textsc{number} and \textsc{case}. Loss of inflection occurs via a complex series of incremental changes, correlated with inflectional class, phonological shape and \textsc{gender}; the loss of \textsc{case} contrasts involves regular sound change, analogy and syntactic change, while the loss of \textsc{number} contrasts is principally due to sound change. While contrast in \textsc{number} is consistently retained longer than contrast in \textsc{case}, it is also noteworthy that, while the overall trend is towards reduction, change is not fast-paced as lexical items with and without given contrasts coexist over several centuries.

---Katja Friedewald studies the history of the French particle \emph{voilà} `here it is/there you go' in her chapter ``French « voilà »: an uninflectable form arising from an inflecting verb''. This particle is uninflecting, a situation that Friedewald argues is quite surprising with regard to its syntactic behavior because the various configurations where it occurs suggest that \emph{voilà} actually	behaves like a verb. In particular, \emph{voilà} takes complements and it assigns accusative case to them, and it can even be embedded as the only element of a relative clause. From a diachronic point of view, the origins of the particle lie in the combination of the imperative form of the verb VEOIR `see' from Old French with the locative adverb \emph{là} `there'. It was hence originally a multiword construction that shows certain similarities to constructional uninflectedness (Spencer, this volume). For Modern French, Friedewald argues that \emph{voilà} cannot be regarded as the singular imperative form of VOIR `see', but rather is an uninflected word since it does not have a \textsc{plural imperative} form, as it did in Medieval French. Corpus data suggest that after a period of parallel use, the \textsc{singular imperative} form is the one usually used from the 16\textsuperscript{th} c. onwards.

---In their chapter, ``Emerging uninflectedness in French clipped verbs'', Maria Copot, Ninoh A. Da Silva, Ahmed Beji, Arno Watiez and Olivier Bonami show that truncation in French nouns (\emph{pneumatique} \textgreater{} \emph{pneu}; \emph{introduction} \textgreater{} \emph{intro}) is well established and well documented, but recent years have seen the emergence of a parallel truncation process in verbs (\emph{je me déconnecte} \textgreater{} \emph{je me déco} `I log off'). The combination of apocope and of the suffixal nature French conjugation provide an interesting new source of uninflectedness. The authors devise a computational method for retrieving truncated verbs from a corpus and analyse the retrieved forms in order to extract generalizations on the phenomenon. They first extract all verbal tokens in the corpus which are a substring of an inflected verb in the French lexicon (e.g. \emph{déco} gets extracted since it is a substring of verbs such as \emph{déco\_nnecter}, \emph{déco\_rer}, \emph{déco\_der}, etc.). The full forms matched by the substring constitute potential matches for the corresponding full form. The best match is chosen by evaluating how well it fits the context, on the basis of methods from distributional semantics. Human evaluation	reveals that the best match is indeed the correct corresponding full form in 88\% of the cases. The conclusion of Copot et al.'s computational corpus study is that the phenomenon is productive, not limited to a subset of lexicalised truncated forms, nor to a restricted set of paradigm cells. The authors further argue that from the point of view of theoretical morphology, it is interesting to note that, while inflectional suffixes are fully lost in truncated verbs, stem allomorphy sometimes leads to the preservation of partial inflectional properties. For instance, among truncated forms of AVOIR `have', a form such as \emph{nous av} from \emph{nous av­\_ons} or \emph{nous av\_ions} `we have'/'we had' signal a present or imperfect, while \emph{j'aur} from \emph{j'aur\_ai} or \emph{j'aur\_ais} `I'll have/I'd have' signals a future or conditional. This raises interesting questions on the division of labour between stem allomorphy and affixal exponence in inflection systems, which the authors also discuss.

---The chapter ``On the emergence of uninflectedness: The case of incipient verbal inflection dropping in present-day French'' by Vincent Renner and Adam Renwick studies the adaptation of verbs from English to French inflectional para\-digms resulting in the recent emergence of uninflected variants for a large number of such verbs (e.g., \emph{je viens de `play' cette vidéo\ldots{}} `I have just played this video\ldots'). The authors' data come from a corpus of 50 million French tweets showing that the lack of inflection on verbs from English is a quantitatively remarkable fact, with thousands of daily attestations in 2021. Similarly, the distribution between inflected and uninflected infinitive forms in a sample of 50 verbs is quite striking as practically all verbs are predominantly uninflected. Renner and Renwick also argue that verbal uninflectedness in present-day French is not limited to borrowings from English. It appears repeatedly at the margins of the lexicon, in colloquial language of slang origin such as verlan and in borrowings from Romani, as well as in the newly emerged French-based contact languages of western Africa Nouchi and the so-called Camfranglais, as verbs, mostly borrowed from English, Cameroonian Pidgin English and Duala do not display inflection, especially in the infinitive and the past participle. The authors claim that while uninflectedness appears to be a marker of limited assimilation, flagging foreignness and peripherality, in the context of social media, uninflectedness seems to be used as a marker of textual genre and in-group membership.

---Anna Thornton and Paolo D'Achille address the development of uninflectedness over time in Italian nouns and adjectives, contributing to a typology of factors that allow them to predict whether an item is uninflectable. Even the factors that correlate with a noun's uninflectability increase and complexify over time: while in the 13\textsuperscript{th} century only nouns ending in a stressed vowel were uninflected, later other classes of nouns become (entirely or partially) uninflectable: nouns ending in unstressed \emph{-e}, in \emph{-i}, in a consonant, and finally \textsc{masculine} nouns in \emph{-a} and \textsc{feminine} nouns in \emph{-o}, where a noun's \textsc{gender} plays a role in determining its uninflectability. These nouns are understood to have developed uninflectability because they contrast with the most system-adequate classes of inflectable nouns in Italian (\textsc{masculine} nouns in \emph{-o} and \textsc{feminine} nouns in \emph{-a}). Thornton and D'Achille also study the behaviour of uninflectable Italian adjectives in diachrony for the first time. Adjectives, having a four-cell paradigm, have a potential for partial uninflectedness within a single paradigm. Preliminary results seem to point to a smaller rate of uninflectedness in Italian adjectives with respect to nouns; this possibly reflects a more pronounced necessity of having overt exponents of contextual inflectional feature values vs. inherent ones.

\medskip

\noindent IV. Uninflectedness by specific word class:

\medskip

---Jerzy Gaszewski's chapter, ``Uninflectedness as a rule in Polish, an inflected language'', examines the phenomenon of uninflected nouns in Polish, a language characterized by rich inflectional morphology involving distinction in \textsc{case}, \textsc{number} and declension classes. Uninflected nouns include stems with phonetic shapes that poorly match existing inflectional patterns (e.g., \emph{kiwi, alibi, Peru}) and nouns denoting culturally distant concepts or foreign names (e.g., \emph{karate, San Francisco, Sukarno}). However, usage can vary, as seen with the inflected noun \emph{papaja} vs uninflected \emph{mango}, despite the fact that both nouns refer to exotic fruits and fit productive declension classes. A particularly regular feature of Polish uninflected nouns is found in animate \textsc{masculine} nouns used for female referents, such as common nouns	(\emph{profesor, minister, architekt}) and most surnames not ending in \emph{-ska/-cka/-dzka}. This uninflected form signals \textsc{feminine}, although it is debated whether these forms and their inflected \textsc{masculine} counterparts have become separate lexemes. These \textsc{feminine} forms have appeared in Polish at least since the 19\textsuperscript{th} century and have competed with normally inflected \textsc{feminine} forms derived with productive suffixes (e.g., \emph{lekarka} for `female medical doctor'). The choice between inflected and uninflected words for female referents has been ideologically charged, with debates in the early 20th c. re-emerging today. The current debate could potentially eliminate the uninflected pattern. However, only common nouns are subject to this controversy, not surnames, suggesting that the development of uninflected \textsc{feminine} forms may be influenced more by sociocultural factors than by linguistic systemic pressures.

---Viktor Köhlich's chapter examines the Japanese word class known as ``rentaishi'' (noun-modifying words), which has been largely overlooked in Western linguistics. Rentaishi words were introduced in the early 20\textsuperscript{th} century to Japanese to align with adjectives in Western fusional languages but have language-specific and typological significance in the context of uninflectedness. Japanese features rich inflectional morphology for verbs and adjectives; in contrast, rentaishi words are characterized by their restriction to attributive positions, inability to inflect, and morphological diversity (e.g., \emph{aru} `a certain'; \emph{rei-no} `mere'), distinguishing them from major word classes based on morphology. Köhlich concludes by discussing whether rentaishi words should remain a distinct word class in Japanese or if these lexemes should be reclassified as uninflecting members of other word classes or even as part of the adjective group they were originally separated from.

\medskip

\noindent V. Diachronic stability of uninflectedness:

\medskip

---In his chapter, ``The diachronic stability of uninflectedness in Berber,'' Lameen Souag argues that comparative data from varieties of Berber, an Afroasiatic language family indigenous to northern Africa, demonstrate that selective uninflectedness can remain stable over a period of approximately two millennia. He starts with the observation that most Berber varieties possess a robust system in which nouns are inflected for a feature known as ``state.'' Typically, this involves two forms: the citation form (referred to as the ``free state'') and a marked form (traditionally misnamed the ``construct state''). The exact distribution of these forms varies across different languages, occasionally interacting with information structure, but they are primarily used for postverbal nominatives and objects of prepositions. Souag shows that in the different Berber varieties, ``state'' inflection applies only to nouns with a gender-marked prefix, which a significant minority of nouns lack. In some instances, even nouns with such a prefix remain uninflected due to phonological reasons. Many (though not all) nouns borrowed from Arabic or Romance languages are uninflected for state, but this also applies to several inherited terms. Some of these terms, such as basic kinship terms, are unambiguously reconstructible as uninflected in proto-Berber. While there are plausible instances where state marking has been extended to previously uninflected nouns, no language has generalized this across the board to make all nouns inflected.

{\sloppy\printbibliography[heading=subbibliography,notkeyword=this]}
\end{document}
